\section*{Chapitre \ref{ch:b2:climate-proxies} --- Archives climatiques et indicateurs}

\begin{solution}{exo:b2:climate-proxies:1}
(a) $\delta = (27.5 \times (-30) + 1350 \times 0)/(1350 + 27.5) = -825/1377.5 = -0.60$ pour mille.
(b) L'océan contiendrait alors toute l'eau : $-0.60$ pour mille, \num{0.60} pour mille plus léger
qu'aujourd'hui. (c) $27.5/1377.5 = \qty{2.0}{\percent}$.
\end{solution}

\begin{solution}{exo:b2:climate-proxies:2}
(a) $T = 16.1 - 4.76\,(2.50 - (0 - 0.27)) = 16.1 - 4.76 \times 2.77 = 16.1 - 13.19 =
\qty{2.9}{\celsius}$. (b) $T = 16.1 - 4.76\,(2.50 - 0.73) = 16.1 - 8.43 = \qty{7.7}{\celsius}$~: la même
coquille donne une température presque cinq degrés plus élevée si l'eau était plus lourde, ce qui montre
qu'il faut connaître la composition de l'eau. (c) La coquille est plus lourde de $2.71 - 2.50 = 0.21$ pour mille,
donc son eau était plus froide de $4.76 \times 0.21 = \qty{1.0}{\celsius}$, à \qty{1.9}{\celsius}.
\end{solution}

\begin{solution}{exo:b2:climate-proxies:3}
(a) $0.4\,\mathrm{e}^{0.9} = \qty{0.98}{mmol/mol}$ et $0.4\,\mathrm{e}^{1.8} =
\qty{2.42}{mmol/mol}$. (b) Le facteur vaut $\mathrm{e}^{0.9} = 2.46$ ; par degré, $\mathrm{e}^{0.09}
= 1.094$, soit une croissance de \qty{9.4}{\percent} par degré. (c) $T = \ln(\text{Mg/Ca}/0.4)/0.09$~:
$\ln 5/0.09 = \qty{17.9}{\celsius}$ et $\ln 2.5/0.09 = \qty{10.2}{\celsius}$.
\end{solution}

\begin{solution}{exo:b2:climate-proxies:4}
(a) $2668 - 2213 + 1 = 456$ ans. (b) Un seul arbre peut présenter localement une absence de cerne, ou un cerne
illisible ; la mise en correspondance des motifs de cernes larges et étroits entre plusieurs arbres révèle les cernes
manquants et fixe l'année de chaque cerne, et l'accord entre arbres mesure la
fiabilité de la datation. (c) À la limite arborée, un seul facteur, le froid de la saison de croissance,
limite tous les arbres, qui partagent donc le même motif de bonnes et de mauvaises années ; dans une
plaine douce et humide, aucun facteur unique ne domine, les motifs des arbres diffèrent, et l'\omterm{def:b2:climate-proxies:dendro}{interdatation} peut échouer.
\end{solution}

\begin{solution}{exo:b2:climate-proxies:5}
(a) $125 \times 40/(3682.2 - 125) = 5000/3557.2 = 1.41$ pour mille ; $134 \times 40/3548.2 =
5360/3548.2 = 1.51$ pour mille. (b) Avec $H - h = 3552.2$~: $3900/3552.2 = 1.10$, $4550/3552.2 = 1.28$
et $6500/3552.2 = 1.83$ pour mille. (c) $\delta_{\mathrm{ice}} = -1.0 \times 3552.2/130 = -27.3$ pour
mille, plus lourd que tout inlandsis actuel (Groenland $-30$ à $-35$, inlandsis nord-américain
$-40$, Antarctique jusqu'à $-50$). Le bilan de masse avec une glace réaliste et la mesure dans les eaux interstitielles
ne concordent pas ; l'écart d'environ un demi pour mille est une incertitude réelle de la méthode.
\end{solution}

\begin{solution}{exo:b2:climate-proxies:6}
(a) $0.372\,\mathrm{e}^{0.0936T} = 0.4\,\mathrm{e}^{0.09T}$ donne $\mathrm{e}^{0.0036T} =
0.4/0.372 = 1.0753$, donc $T = \ln 1.0753/0.0036 = \qty{20.2}{\celsius}$. (b) À
\qty{2}{\celsius}~: $0.4\,\mathrm{e}^{0.18} = 0.479$ et $0.372\,\mathrm{e}^{0.1872} = 0.449$, donc le rapport
ajusté vaut $0.449/0.479 = 0.937$, soit \qty{6.3}{\percent} de moins. (c) $T = \ln(0.449/0.4)/0.09 =
\ln 1.1225/0.09 = \qty{1.3}{\celsius}$~: la loi prévue lit la coquille environ
\qty{0.7}{\celsius} trop froide. À basse température, le choix de la calibration compte au niveau d'un
degré.
\end{solution}

\begin{solution}{exo:b2:climate-proxies:7}
(a) De l'Holocène à GI-1, \qty{112.19}{m} de glace pour 2989 ans~: $11\,219/2989 =
\qty{3.75}{cm}$ par an ; de GI-16 à GI-17, \qty{17.89}{m} pour 1160 ans~: $1789/1160 =
\qty{1.54}{cm}$ par an. (b) La glace
s'écoule~: à mesure qu'une couche est enfouie, elle s'étale latéralement et s'amincit, et les couches plus profondes et plus anciennes
ont été amincies plus longtemps. (c) Lorsque les couches deviennent trop minces, un centimètre ou moins, les instruments
ne peuvent plus distinguer une année de la suivante, et le \omterm{def:b2:climate-proxies:age-model}{modèle d'âge} doit alors recourir à des points de calage ou à des modèles d'écoulement
à la place.
\end{solution}

\begin{solution}{exo:b2:climate-proxies:8}
(a) La fraction de l'intervalle vaut $(1550 - 1492.45)/(1604.64 - 1492.45) = 57.55/112.19 = 0.513$,
donc l'âge vaut $11\,703 + 0.513 \times 2989 = 11\,703 + 1533 = \num{13236}$ ans b2k. (b) 1950 est 50
ans avant 2000 : \num{13186} ans avant 1950. (c) $50 + 0.513 \times (93 - 50) = 72$ ans.
\end{solution}

\begin{solution}{exo:b2:climate-proxies:9}
(a) Chaque couche incertaine apporte une demi-année : $2 \times 93 = 186$ et $2 \times 1287 = 2574$
couches incertaines. (b) $\frac12\sqrt{186} = 6.8$ ans et $\frac12\sqrt{2574} = 25.4$ ans.
(c) $93/6.8 = 14$ et $1287/25.4 = 51$. Le rapport de $\frac12 n$ à $\frac12\sqrt{n}$ vaut
$\sqrt{n}$ : il croît comme la racine carrée du nombre de couches incertaines, si bien que le choix entre les deux règles
compte de plus en plus avec la profondeur.
\end{solution}

\begin{solution}{exo:b2:climate-proxies:10}
(a) Sur la \omterm{prop:b2:isotope-geochemistry:mwl}{droite des eaux météoriques}, $\delta$D varie huit fois plus vite que $\delta^{18}$O, donc la
pente vaut $8 \times 0.8 = 6.4$ pour mille de $\delta$D par degré. La baisse vaut $441.7 - 396.6 = 45.1$
pour mille, d'où un refroidissement de $45.1/6.4 = \qty{7.0}{\celsius}$. (b) L'océan glaciaire était
plus lourd d'environ $8 \times 1.0 = 8$ pour mille en $\delta$D ; la neige formée à partir de l'océan actuel aurait été
plus légère d'autant, soit $-441.7 - 8 = -449.7$ pour mille. La baisse devient $53.1$ pour mille et le refroidissement
$53.1/6.4 = \qty{8.3}{\celsius}$. (c) $9.39 - 0.09 = \qty{9.30}{\celsius}$, donc $45.1/9.30 = 4.85$
pour mille par degré. La série publiée est aussi corrigée des changements d'altitude de l'inlandsis ;
sa relation de température provient de modèles qui simulent les isotopes de l'eau,
et non de la seule pente spatiale.
\end{solution}

\begin{solution}{exo:b2:climate-proxies:11}
(a) Les logarithmes des rapports valent $-0.528$, $-0.223$, $0.191$, $0.501$, $0.904$, $1.227$ ; leur
moyenne vaut $0.345$, et la température moyenne vaut 14. Avec $S_{TT} = \sum (T - 14)^2 = 280$ et $S_{T\ell} =
\sum (T - 14)(\ell - 0.345) = 24.93$, $A = 24.93/280 = \qty{0.0890}{\per\celsius}$ et $\ln B =
0.345 - 0.0890 \times 14 = -0.901$, donc $B = \qty{0.406}{mmol/mol}$. (b) Les résidus sont
$+0.017$, $-0.034$, $+0.023$, $-0.023$, $+0.025$ et $-0.009$, soit quelques pour cent du rapport, sans
tendance ; l'ajustement est proche de $0.4\,\mathrm{e}^{0.09T}$. (c) $T = (\ln 1.30 + 0.901)/0.0890 = (0.262 +
0.901)/0.0890 = \qty{13.1}{\celsius}$. (d) $T = (\ln 0.45 + 0.901)/0.0890 = (-0.799 +
0.901)/0.0890 = \qty{1.2}{\celsius}$. Cette valeur se situe sous l'échantillon le plus froid de la calibration,
\qty{4}{\celsius} : la loi est extrapolée, et, aux basses températures, le Mg/Ca des coquilles est connu pour s'écarter
de la loi exponentielle sous l'influence de l'ion carbonate.
\end{solution}

\begin{solution}{exo:b2:climate-proxies:12}
(a) Par unité de surface océanique, l'eau de l'océan actuel, $H$ à 0, se partage entre les réserves
$h_i$ à $\delta_i$ et l'océan restant $H - h$, avec $h = \sum h_i$~: $0 = \sum h_i\delta_i +
(H - h)\,\Delta\delta$. Après la première étape, $\Delta\delta = 60 \times 30/(3682.2 - 60) =
1800/3622.2 = 0.50$ pour mille ; après la seconde, $(1800 + 70 \times 45)/(3682.2 - 130) =
4950/3552.2 = 1.39$ pour mille. (b) $-4950/130 = -38.1$ pour mille. (c) La règle est $1.39/130 = 0.0107$
pour mille par mètre, et donne $0.50/0.0107 = \qty{47}{m}$ pour la première étape au lieu de
\qty{60}{m} : la première glace était plus lourde que la moyenne, si bien que chacun de ses mètres décalait l'océan
moins. Une règle linéaire calibrée sur un seul état se trompe sur les autres lorsque la composition moyenne de la
glace change.
\end{solution}

\begin{solution}{pb:b2:climate-proxies:1}
\textbf{1.} $361.84 \times 10^6\ \unit{km^2} \times \qty{0.130}{km} = \qty{4.70e7}{km^3}$, soit environ 47
millions de kilomètres cubes.
\textbf{2.} $130 \times 40/(3682.2 - 130) = 5200/3552.2 = 1.46$ pour mille.
\textbf{3.} $5000/3557.2 = 1.41$ et $5360/3548.2 = 1.51$ pour mille.
\textbf{4.} $4550/3552.2 = 1.28$ et $6500/3552.2 = 1.83$ pour mille. La composition de la glace
modifie la réponse de \num{0.55} pour mille, le niveau de la mer de \num{0.10} : la glace compte davantage.
\textbf{5.} $\delta_{\mathrm{ice}} = -1.0 \times 3552.2/130 = -27.3$ pour mille.
\textbf{6.} $40/3682.2 = 0.0109$ pour mille par mètre, ce qui s'arrondit au \num{0.011} des niveaux marins
datés.
\textbf{7.} À $\delta^{18}\mathrm{O}_{\mathrm{sw}}$ constant, $\mathrm{d}T = -4.76\,\mathrm{d}\delta^{18}
\mathrm{O}_{\mathrm{b}}$, donc $\mathrm{d}T = -1$ donne $\mathrm{d}\delta^{18}\mathrm{O}_{\mathrm{b}} =
1/4.76 = 0.21$ pour mille.
\textbf{8.} $\delta^{18}\mathrm{O}_{\mathrm{b}} = (16.1 - T)/4.76 + \delta^{18}\mathrm{O}_{\mathrm{sw}}
- 0.27 = 14.1/4.76 - 0.27 = 2.96 - 0.27 = 2.69$ pour mille.
\textbf{9.} $17.1/4.76 + 1.0 - 0.27 = 3.59 + 0.73 = 4.32$ pour mille. La variation vaut $4.32 - 2.69 =
1.63$ pour mille : \num{1.00} dû à l'eau et $3.0 \times 0.21 = 0.63$ dû au refroidissement de
\qty{3}{\celsius}.
\textbf{10.} $T = \ln(2.0/0.4)/0.09 = \ln 5/0.09 = \qty{17.9}{\celsius}$.
\textbf{11.} $0.18/0.1 = \qty{1.8}{\celsius}$ de refroidissement.
\textbf{12.} Le Mg/Ca est un rapport de deux éléments dans la calcite, fixé par la température de la
calcification ; il ne fait pas intervenir les isotopes de l'oxygène, et il ignore donc la composition de
l'eau que la glace modifie.
\textbf{13.} $4.92 - 3.23 = 1.69$ pour mille.
\textbf{14.} En résolvant $\Delta = h \times 40/(H - h)$ en $h$ : $h = H\Delta/(\Delta + 40) = 3682.2
\times 1.69/41.69 = \qty{149}{m}$, plus que les 125 à \qty{134}{m} de l'indice de niveau de la mer :
une partie de la variation doit être due au refroidissement.
\textbf{15.} $1.69 - 1.46 = 0.23$ pour mille, et $4.76 \times 0.23 = \qty{1.1}{\celsius}$.
\textbf{16.} $1.69 - 1.00 = 0.69$ pour mille, et $4.76 \times 0.69 = \qty{3.3}{\celsius}$ (3.28).
\textbf{17.} $0.23/0.25 = \qty{0.9}{\celsius}$ et $0.69/0.25 = \qty{2.8}{\celsius}$.
\textbf{18.} L'eau profonde nord-atlantique est passée de 2 à \qty{4}{\celsius} à \qty{-1.1}{\celsius},
de 3.1 à \qty{5.1}{\celsius} ; l'eau de fond antarctique est passée de 0 à \qty{1}{\celsius} à
\qty{-1.9}{\celsius}, de 1.9 à \qty{2.9}{\celsius} ; le Mg/Ca de l'Atlantique Nord-Ouest donne
\qty{1.8}{\celsius}. Toutes se situent entre environ 2 et \qty{5}{\celsius}, autour des 2.8 à
\qty{3.3}{\celsius} de la correction par les eaux interstitielles, et au-dessus du \qty{1}{\celsius} du bilan de
masse avec une glace à $-40$ pour mille.
\textbf{19.} $-396.6 - (-441.7) = 45.1$ pour mille.
\textbf{20.} $8 \times 0.8 = 6.4$ pour mille par degré, donc $45.1/6.4 = \qty{7.0}{\celsius}$.
\textbf{21.} $9.39 - 0.09 = \qty{9.30}{\celsius}$, et $45.1/9.30 = 4.85$ pour mille par degré.
\textbf{22.} Dome C, \qty{9.3}{\celsius}, plus que la surface du globe, de 5 à \qty{7}{\celsius}, plus
que l'océan profond, d'environ \qty{3}{\celsius}.
\textbf{23.} L'eau profonde est déjà froide aujourd'hui, de 0 à \qty{4}{\celsius}, et l'eau de mer ne peut guère se refroidir
au-dessous de son point de congélation, que l'eau de fond glaciaire a presque atteint à
\qty{-1.9}{\celsius} : l'océan profond n'avait que quelques degrés à perdre, alors que l'air au-dessus de l'Antarctique
n'a pas un tel plancher.
\textbf{24.} Avec la correction par les eaux interstitielles, l'océan profond était environ \textbf{\qty{3.3}{\celsius}
plus froid} au dernier maximum glaciaire qu'aujourd'hui, proche des 2 à \qty{5}{\celsius} des estimations
indépendantes pour l'eau profonde.
\end{solution}
